\section{The group and a finite-dimensional representation}
\label{sec:group}

We construct a finite-dimensional representation inequivalent to each of
its positive pullbacks along a fixed group endomorphism. The same group
configuration appears in the nonfiniteness arguments of
Sauers~\cite[Section~2.4]{Sauers2026NonMF} and
Eckhardt~\cite[Theorem~3.1 and Example~3.2]{Eckhardt2026}.
De Cornulier~\cite[Proposition~1.1 and its proof]{deCornulier2007}
records the property~(T) affine group and its finite-index scalar
self-embeddings. We use the dyadic case.

We use column vectors.  In the group $G$ of Theorem~\ref{thm:main},
write elements as $(v,M,k)$, so that
\[
 (v,M,k)(w,N,\ell)=(v+2^kMw,MN,k+\ell).
\]
Set
\[
 \Gamma=\{(v,M,0):v\in\Z^3,\ M\in\SL_3(\Z)\},
 \qquad t=(0,I,1).
\]
We suppress the final coordinate when writing elements of $\Gamma$, and
define
\[
 \theta:\Gamma\longrightarrow\Gamma,\qquad
 \theta(v,M)=(2v,M),
 \qquad \Lambda=\theta(\Gamma).
\]

\begin{lemma}\label{lem:group-data}
The group $\Gamma$ is finitely generated and has property~(T).  The map
$\theta$ is injective,
\[
 [\Gamma:\Lambda]=8,
 \qquad t\gamma t^{-1}=\theta(\gamma)\quad(\gamma\in\Gamma).
\]
There is a nonzero finite-dimensional irreducible unitary representation
$\pi_*$ of $\Gamma$ that factors through reduction modulo~$2$ and is
nontrivial on the translation subgroup $\Z^3$.  In particular,
\begin{equation}\label{eq:group-aperiodic}
 \pi_*\not\simeq\pi_*\circ\theta^j\qquad(j\geq1).
\end{equation}
\end{lemma}

\begin{proof}
Property~(T) is the case $n=3$, $m=1$ of
\cite[Proposition~1.1]{deCornulier2007}, with $\SL_3(\Z)$ acting on one
copy of its standard rank-three integral module.  For finite generation,
take the three standard translations and the elementary matrices
$I+E_{ij}$, $i\ne j$, where $E_{ij}$ is the matrix unit.  These matrices
generate $\SL_3(\Z)$ by integer row reduction.

The group law shows that $\theta$ is an injective homomorphism and that
conjugation by $t$ implements it.  Its image is
$\Lambda=2\Z^3\rtimes\SL_3(\Z)$.  The elements
\[
 r_a=(a,I),\qquad a\in\{0,1\}^3,
\]
are representatives of its eight left cosets: writing $v=a+2w$ gives
$(v,M)=r_a(2w,M)$, and two such representatives give the same coset
only when they have the same residue modulo~$2$.

Let $F$ be the finite image of the reduction homomorphism
\[
 \Gamma\longrightarrow (\Z/2\Z)^3\rtimes\SL_3(\Z/2\Z),
 \qquad (v,M)\longmapsto(v\bmod2,M\bmod2).
\]
The translation subgroup $(\Z/2\Z)^3$ acts nontrivially in the left
regular representation of $F$.  Decomposing this representation into
irreducibles therefore gives an irreducible constituent that is
nontrivial on translations.  Its pullback to $\Gamma$ is the required
$\pi_*$.  Indeed, $\pi_*$ factors through reduction modulo~$2$, whereas
$\theta^j(v,I)=(2^jv,I)$.  Thus $\pi_*\circ\theta^j$ is trivial on all
translations for every $j\geq1$, and $\pi_*$ is not.  This proves
\eqref{eq:group-aperiodic}.
\end{proof}

The group $G$ is countable, so its full group $C^*$-algebra $C^*(G)$ is
separable and unital.  Every unitary representation of $G$ extends to
$C^*(G)$.  We will use this universal property after inducing $\pi_*$
from $\Gamma$ to $G$.
